% ====================================================================
% ФАЙЛ: tasks/01_mechanics/energy_conservation_bank_part1.tex
% ТЕМА: КИНЕТИЧЕСКАЯ И ПОТЕНЦИАЛЬНАЯ ЭНЕРГИЯ. ЗСЭ
% ПАЧКА 1: ВАРИАНТЫ 1–10
% ====================================================================

% === ЗАДАЧА 1 ===
% Тег: #МЕХАНИКА #КИМ_06 #ВАРИАНТ_01
\begin{taskbasic}[code={\label{task:energy_v1_n6}}]
\textbf{Задача 1.} Груз, привязанный к длинной нерастяжимой нити, отклонили от положения равновесия на малый угол и в момент времени $t=0$ отпустили с нулевой начальной скоростью (см. рисунок). На графиках А и Б показано изменение физических величин, характеризующих движение груза после этого. $T$ — период колебаний. Сопротивлением воздуха пренебречь. Потенциальная энергия груза отсчитывается от положения равновесия.

\nopagebreak\smallskip
\begin{center}
\begin{stylebasic}
\begin{tabular}{cc}
\begin{tikzpicture}[scale=0.7, every node/.style={font=\footnotesize}]
    \draw[thick] (-0.8,2.5) -- (0.8,2.5);
    \foreach \x in {-0.7,-0.5,-0.3,-0.1,0.1,0.3,0.5,0.7}
        \draw (\x,2.5) -- (\x+0.1,2.7);
    \draw[dashed] (0,2.5) -- (0,0);
    \draw[thick] (0,2.5) -- (-0.8,0.5);
    \draw[thick, fill=gray!40] (-0.8,0.5) rectangle (-0.4,1.1);
    \draw[-{Stealth[scale=0.7]}] (-1.2,-0.2) -- (1.5,-0.2) node[right] {$x$};
    \node[below] at (0,-0.2) {0};
\end{tikzpicture} 
&
\begin{tabular}{l}
Установите соответствие между графиками и физическими\\
величинами, зависимость которых от времени эти графики\\
могут представлять.\\
\end{tabular}
\end{tabular}
\end{stylebasic}
\end{center}

К каждой позиции первого столбца подберите соответствующую позицию из второго столбца и запишите в таблицу выбранные цифры под соответствующими буквами.

\nopagebreak\smallskip
\begin{center}
\begin{tabular}{cc}
\textbf{ГРАФИКИ} & \textbf{ФИЗИЧЕСКИЕ ВЕЛИЧИНЫ} \\
\begin{tikzpicture}[scale=0.55, every node/.style={font=\tiny}]
    \draw[-{Stealth[scale=0.7]}] (0,0) -- (4.5,0) node[right] {$t$};
    \draw[-{Stealth[scale=0.7]}] (0,-1.5) -- (0,2) node[above] {А};
    \draw[ultra thick, brandPrimary] plot[domain=0:4, samples=100] (\x, {1.2*cos(\x*180/2)});
    \draw[dashed] (4,0) -- (4,1.2);
    \node[below] at (4,0) {$T$};
    \node[below left] at (0,0) {0};
\end{tikzpicture}
&
\begin{tabular}{l}
1) проекция скорости $v_x$ \\
2) координата $x$ \\
3) проекция ускорения $a_x$ \\
4) потенциальная энергия груза $E_{\text{п}}$
\end{tabular} \\
\begin{tikzpicture}[scale=0.55, every node/.style={font=\tiny}]
    \draw[-{Stealth[scale=0.7]}] (0,0) -- (4.5,0) node[right] {$t$};
    \draw[-{Stealth[scale=0.7]}] (0,-1.5) -- (0,1.5) node[above] {Б};
    \draw[ultra thick, brandPrimary] plot[domain=0:4, samples=100] (\x, {1.0*sin(\x*180/2)});
    \draw[dashed] (4,0) -- (4,0);
    \node[below] at (4,0) {$T$};
    \node[below left] at (0,0) {0};
\end{tikzpicture} & 
\end{tabular}
\end{center}

\begin{center}
\begin{tabular}{|c|c|}
\hline
А & Б \\ \hline
 & \\ \hline
\end{tabular}
\end{center}

\kim{6}
\end{taskbasic}
\answer{41}

% === ЗАДАЧА 2 ===
% Тег: #МЕХАНИКА #КИМ_06 #ВАРИАНТ_02
\begin{taskbasic}[code={\label{task:energy_v2_n6}}]
\textbf{Задача 2.} Груз, привязанный к длинной нерастяжимой нити, отклонили от положения равновесия на малый угол и в момент времени $t=0$ отпустили с нулевой начальной скоростью (см. рисунок). На графиках А и Б показано изменение физических величин, характеризующих движение груза после этого. $T$ — период колебаний. Сопротивлением воздуха пренебречь. Потенциальная энергия груза отсчитывается от положения равновесия.

\nopagebreak\smallskip
\begin{center}
\begin{stylebasic}
\begin{tabular}{cc}
\begin{tikzpicture}[scale=0.7, every node/.style={font=\footnotesize}]
    \draw[thick] (-0.8,2.5) -- (0.8,2.5);
    \foreach \x in {-0.7,-0.5,-0.3,-0.1,0.1,0.3,0.5,0.7}
        \draw (\x,2.5) -- (\x+0.1,2.7);
    \draw[dashed] (0,2.5) -- (0,0);
    \draw[thick] (0,2.5) -- (-0.8,0.5);
    \draw[thick, fill=gray!40] (-0.8,0.5) rectangle (-0.4,1.1);
    \draw[-{Stealth[scale=0.7]}] (-1.2,-0.2) -- (1.5,-0.2) node[right] {$x$};
    \node[below] at (0,-0.2) {0};
\end{tikzpicture} 
&
\begin{tabular}{l}
Установите соответствие между графиками и физическими\\
величинами, зависимость которых от времени эти графики\\
могут представлять.\\
\end{tabular}
\end{tabular}
\end{stylebasic}
\end{center}

К каждой позиции первого столбца подберите соответствующую позицию из второго столбца и запишите в таблицу выбранные цифры под соответствующими буквами.

\nopagebreak\smallskip
\begin{center}
\begin{tabular}{cc}
\textbf{ГРАФИКИ} & \textbf{ФИЗИЧЕСКИЕ ВЕЛИЧИНЫ} \\
\begin{tikzpicture}[scale=0.55, every node/.style={font=\tiny}]
    \draw[-{Stealth[scale=0.7]}] (0,0) -- (4.5,0) node[right] {$t$};
    \draw[-{Stealth[scale=0.7]}] (0,-1.5) -- (0,2) node[above] {А};
    \draw[ultra thick, brandPrimary] plot[domain=0:4, samples=100] (\x, {1.2*cos(\x*180/2)});
    \draw[dashed] (4,0) -- (4,1.2);
    \node[below] at (4,0) {$T$};
    \node[below left] at (0,0) {0};
\end{tikzpicture}
&
\begin{tabular}{l}
1) проекция скорости $v_x$ \\
2) координата $x$ \\
3) проекция ускорения $a_x$ \\
4) потенциальная энергия груза $E_{\text{п}}$
\end{tabular} \\
\begin{tikzpicture}[scale=0.55, every node/.style={font=\tiny}]
    \draw[-{Stealth[scale=0.7]}] (0,0) -- (4.5,0) node[right] {$t$};
    \draw[-{Stealth[scale=0.7]}] (0,-1.5) -- (0,1.5) node[above] {Б};
    \draw[ultra thick, brandPrimary] plot[domain=0:4, samples=100] (\x, {-1.0*cos(\x*180/2)});
    \draw[dashed] (4,0) -- (4,0);
    \node[below] at (4,0) {$T$};
    \node[below left] at (0,0) {0};
\end{tikzpicture} & 
\end{tabular}
\end{center}

\begin{center}
\begin{tabular}{|c|c|}
\hline
А & Б \\ \hline
 & \\ \hline
\end{tabular}
\end{center}

\kim{6}
\end{taskbasic}
\answer{42}

% === ЗАДАЧА 3 ===
% Тег: #МЕХАНИКА #КИМ_03 #ВАРИАНТ_07
\begin{taskbasic}[code={\label{task:energy_v7_n3}}]
\textbf{Задача 3.} У основания гладкой наклонной плоскости шайба обладает скоростью, модуль которой равен $4$~м/с. Определите массу шайбы, если максимальная потенциальная энергия, которую она приобретает при подъёме по плоскости относительно её основания, составляет $0{,}4$~Дж. Сопротивлением воздуха пренебречь.

\kim{3}
\end{taskbasic}
\answer{0{,}05}

% === ЗАДАЧА 4 ===
% Тег: #МЕХАНИКА #КИМ_03 #ВАРИАНТ_08
\begin{taskbasic}[code={\label{task:energy_v8_n3}}]
\textbf{Задача 4.} У основания гладкой наклонной плоскости шайба массой $120$~г обладает скоростью, модуль которой равен $5$~м/с. Определите максимальную потенциальную энергию, которую приобретёт шайба при подъёме по плоскости относительно её основания. Сопротивлением воздуха пренебречь.

\kim{3}
\end{taskbasic}
\answer{1{,}5}